% =====================================================================
%  blueprint/Piezothermoelastic_AntiReflection_Master_Blueprint_V5.tex
%
%  Faithful LaTeX transcription of the V5.0 master blueprint
%  (PDF: Piezothermoelastic_AntiReflection_Master_Blueprint_V5.pdf,
%   13 pages, 4 November 2026).
%
%  TRANSCRIPTION RULE USED: the scientific content is reproduced without
%  alteration. Nothing has been added, removed, reordered or reinterpreted.
%  Where the PDF's two-column status table needed a different LaTeX device
%  (tabular), the words are unchanged. Section numbering, hypothesis labels
%  (H1, H2), falsification rules (F1-F3), equations (1)-(10), gates G0-G7,
%  phases P0-P5 and the Appendix A change log all follow the original.
%
%  COMPILATION STATUS: NO TeX ENGINE IS INSTALLED in this environment
%  (pdflatex, latex, xelatex, lualatex, tectonic and latexmk are all
%  absent). This file has been inspected as text only. It has NOT been
%  compiled and no claim of successful typesetting is made.
% =====================================================================
\documentclass[11pt]{article}
\usepackage{amsmath,amssymb,amsthm,bm}
\usepackage{geometry}
\usepackage{booktabs}
\usepackage{longtable}
\usepackage{array}
\geometry{margin=1in}
\newcommand{\R}{\mathcal{R}}
\newcommand{\T}{\mathcal{T}}
\newcommand{\A}{\mathcal{A}}
\newcommand{\jj}{[\![}
\newcommand{\kk}{]\!]}

\title{Blueprint: Anti-Reflection and Reflection Suppression of Coupled
Waves in Piezothermoelastic Media with Generalized Imperfect Interface
Conditions\\[6pt]
\large Finite Matching Interphase as Primary Candidate; Two-Half-Space
Imperfect Interface as Comparison Baseline}
\author{Research-design document for independent scientific review}
\date{Prepared 4 November 2026 \quad$\vert$\quad Final working edition (V5.0)}

\begin{document}
\maketitle

\begin{center}
\fbox{\parbox{0.92\textwidth}{\centering\textbf{WORKING RESEARCH BLUEPRINT ---
FLEXIBLE; NOT A GUARANTEED RESULT, NOT A MANUSCRIPT, NOT AN APPROVED NOVELTY
CLAIM}}}
\end{center}

\begin{abstract}
The research direction is fixed: reflection suppression of coupled waves in
piezothermoelastic media through generalized imperfect interface conditions.
This document defines one bounded theoretical wave-dynamics paper. The primary
candidate configuration is two dissimilar piezothermoelastic half-spaces
separated by a finite piezothermoelastic matching interphase; the baseline for
comparison is the two-half-space imperfect (zero-thickness) interface. Section 5
now contains a derived scalar proposition showing exactly when a passive
zero-thickness contact can suppress reflection and what it costs in transmitted
energy. Hypothesis H1 is reformulated accordingly, benchmark reproduction of any
single published paper is conditional on a model-compatibility check, and
Lord--Shulman energy flux, interface dissipation and numerical energy balance
are compulsory validation gates. No coupled-field numerical result is reported in
this document. The supervisor controls the model and approval decisions; external
AI agents are restricted to assigned verification and implementation tasks. The
original V4.2 scientific framework is retained, with formal work allocation and
additional closure checks. A historical change log from V4.1 to V4.2 appears in
Appendix A.

\medskip
\noindent\textbf{Status labels.} \textsc{Derived}: derived in this document by
hand and checked where stated. \textsc{Executed}: a computation actually run
(inputs listed). \textsc{Hypothesis}: untested claim. \textsc{Open}: not yet
known or not yet established. Labels apply only to the item they follow.
\end{abstract}

\section{Scientific identity, scope and flexible decisions}

\textbf{Fixed scope.} Coupled piezothermoelastic PDEs; normal-mode analysis;
characteristic determinants; generalized spring/contact boundary conditions;
reflection and transmission coefficients; energy partition; analytical and
numerical validation. Do not add metamaterial engineering, active control,
machine learning, topological physics, or device-design methods.
Gurtin--Murdoch surface elasticity and Klein--Gordon/strain-gradient nonlocality
are \emph{not} included unless a specific demonstrated need arises and is recorded
in the decision log (Section 11).

\textbf{One-paper rule.} One incident wave family, one reference material
pairing, one coherent central contribution. Variants are validation limits or
fallback configurations, not additional papers.

\begin{center}
\begin{tabular}{@{}p{0.16\textwidth}p{0.50\textwidth}p{0.26\textwidth}@{}}
\toprule
\textbf{Element} & \textbf{Current choice} & \textbf{Status} \\
\midrule
Bulk material & Homogeneous, dissimilar, transversely isotropic
  piezothermoelastic half-spaces; plane strain; poling axis declared in the
  sagittal plane & Core \\
Coupled fields & $u_x, u_z, \varphi, \vartheta$ (four scalar variables; three
  physical categories) & Core \\
Thermal theory & Lord--Shulman (LS), pyroelectric coupling constitutively
  consistent & Initial model \\
Primary configuration & Finite piezothermoelastic matching interphase
  $0<z<h$ with passive generalized contact on both faces & Primary candidate \\
Comparison baseline & Two half-spaces with a single zero-thickness generalized
  imperfect interface & Baseline (also the $h\to0$ comparison) \\
Contact law & Passive diagonal spring--dashpot (mechanical) with declared
  electrical and thermal contact laws & Primary mechanism \\
Gurtin--Murdoch; Klein--Gordon/NSGT & Not used unless a demonstrated need
  arises & Excluded by default \\
\bottomrule
\end{tabular}
\end{center}

\section{Closest prior art, research question and hypotheses}

The literature already contains imperfect interfaces in piezothermoelastic
half-spaces, several thermal theories, energy ratios, and finite
thermo-electro-elastic layers. Neither the four coupled scalar variables, an
arbitrary boundary-condition list, nor the scalar quarter-wave rule is original.

\begin{itemize}
  \item Guha, Singh and Das, imperfect interfaces between two dissimilar
    piezothermoelastic half-spaces, \texttt{10.1080/17455030.2019.1610198}
    (Crossref record checked; full text not read).
  \item Karmakar, Sahu and Goyal, pre-stressed functionally graded
    piezothermoelastic interfaces with mechanical/electrical/thermal loose
    bonding, \texttt{10.1177/10812865221099502} (Crossref abstract read; full
    text not read).
  \item Barak, Kumar, Kumar and Gupta, energy analysis at an elastic/nonlocal
    piezothermoelastic interface, \texttt{10.1007/s12648-022-02568-w} (publisher
    page and abstract read; full text behind a subscription, not read). The
    abstract states two reflected and four transmitted waves.
  \item Gupta, Kumar, Kumar and Barak, energy analysis at a nonlocal
    thermoelastic/piezothermoelastic interface,
    \texttt{10.1108/HFF-11-2022-0654} (publisher page and abstract read).
  \item Wang et al., sandwiched thermo-electro-elastic nanoplates with
    dissipation, \texttt{10.1016/j.compstruct.2023.117564} (plates, not
    half-spaces).
  \item Ghazaryan, Piliposyan and Piliposian, shear waves through a periodically
    stratified piezoelectric reflector with imperfect interfaces,
    \texttt{10.1016/j.ijsolstr.2025.113514} (Crossref record checked: Int.\ J.\
    Solids Struct.\ 321, 2025; shear-wave layered competitor).
\end{itemize}

\textbf{Candidate gap, not a novelty claim.} An explicit, physically admissible
condition for coupled-mode reflection suppression, together with a quantitative
account of what the suppression costs in transmitted power, may not be fully
addressed by these studies. This must be checked against the full equations and
figures of each competitor, not titles or abstracts (gate G0).

\textbf{Main question.} For a finite piezothermoelastic interphase with passive
generalized contact on its faces, under what contact, layer and material choices
is the reflected response of an incident quasi-longitudinal or quasi-shear wave
reduced relative to a defined baseline, and how much of the reduction appears as
increased transmission rather than increased dissipation?

\textbf{Definitions} (all fixed before any search). $R$, $T$, $A$ are
incident-power-normalized reflected, transmitted and dissipated fractions defined
in Section 5. The \emph{welded baseline} (subscript $w$) is the same two
half-spaces with perfect mechanical, electrical and thermal contact.
$\Delta R = R - R_w$ and $\Delta T = T - T_w$.

\subsection*{H1 (revised, zero-thickness contact; falsifiable)}

\textbf{Statement.} For some sourced material pair and some passive, diagonal
zero-thickness generalized contact within declared parameter bounds, there exist
$(\omega,\theta)$ such that $\Delta R \le -\varepsilon_R$ (a pre-declared
reflection reduction) \emph{with} a transmission sacrifice no larger than a
pre-declared budget, $\Delta T \ge -\delta_T$, where $\delta_T>0$ is fixed before
the search. H1 does \emph{not} assert $\Delta T \ge 0$.

\textbf{Known counterexample (derived, Section 5).} In the scalar,
normal-incidence, real-impedance model, every passive zero-thickness contact of
either topology satisfies $T\le T_w$ with equality only when the contact vanishes
(Proposition 1(c)). Hence H1 with $\delta_T=0$ is false in that model; a purely
reactive contact has $R\ge R_w$; a dissipative contact reaching $R=0$ has
$T=Z_{\min}/Z_{\max}<T_w$ whenever $Z_1\ne Z_2$.

\textbf{Conditions under which H1 could still hold in the coupled problem (all
\textsc{open}).} Oblique incidence with mode conversion and critical-angle
structure; thermo-electric coupling providing channels that the scalar model
lacks; electrical or thermal contact laws that alter modal impedances. None is
established.

\textbf{Falsification and reporting rules.} (F1) Any search outcome with
$\Delta R\le-\varepsilon_R$ achieved only with $\Delta T<-\delta_T$ is reported as
\emph{dissipation-driven suppression}, not anti-reflection with useful
transmission. (F2) If no point in the declared parameter box meets both
thresholds for a material pair, H1 is rejected for that pair and the negative
result is reported. (F3) $\varepsilon_R$, $\delta_T$ and the parameter box are
recorded before the search and changed only through a logged decision
(Section 11).

\subsection*{H2 (primary configuration; falsifiable)}

Adding a finite interphase of thickness $h$ provides an additional independent
phase parameter and may reach $\Delta R\le-\varepsilon_R$ together with
$\Delta T\ge0$ for sourced parameters, which H1 cannot do in the scalar model.
Scalar motivation: a lossless layer can achieve $R=0$ with $T=1>T_w$
(Section 5.4); the coupled problem differs because the thermal mode is
dissipative. H2 is tested against the H1 baseline with a comparable parameter
budget; it is rejected if no sourced configuration meets
$\Delta R\le-\varepsilon_R$ and $\Delta T\ge0$.

\section{Bulk problem and mathematical closure}

Medium 1 occupies $z<0$ and medium 2 occupies $z>h$ (or $z>0$ for the
zero-thickness baseline); $x$ is tangential. Use 6mm-type piezoelectric symmetry
with a declared poling direction in the sagittal $x$--$z$ plane, small
perturbations about $T_0>0$, and $e^{-i\omega t}$ time dependence. The unknowns
are
\begin{equation}
  \bm{y} = (u_x, u_z, \varphi, \vartheta)^{\mathsf T}, \qquad
  E_i = -\varphi_{,i}, \qquad
  \varepsilon_{ij} = \tfrac12(u_{i,j}+u_{j,i}).
  \tag{1}
\end{equation}

A conventional candidate linear constitutive set (signs and energy potential to be
verified together) is
\begin{align}
  \sigma_{ij} &= C_{ijkl}\varepsilon_{kl} - e_{kij}E_k - \beta_{ij}\vartheta,
  \tag{2}\\
  D_i &= e_{ikl}\varepsilon_{kl} + \epsilon_{ik}E_k + p_i\vartheta, \tag{3}\\
  \sigma_{ij,j} &= \rho\ddot u_i, \qquad D_{i,i} = 0 . \tag{4}
\end{align}

The LS heat equation must be derived from the same potential. A candidate form
requiring independent sign and entropy verification is
\begin{equation}
  \kappa_{ij}\vartheta_{,ij} = (1+\tau_q\partial_t)\Big[\rho c_E\dot\vartheta
  + T_0\big(\beta_{ij}\dot\varepsilon_{ij} + p_i\dot E_i\big)\Big],
  \qquad \bm{q} + \tau_q\dot{\bm{q}} = -\kappa\nabla\vartheta . \tag{5}
\end{equation}
A hand dimensional check shows that each term inside the bracket has units
$\mathrm{W\,m^{-3}}$ (\textsc{derived}).

\textbf{Required restrictions:} material stability, $\rho, c_E, T_0, \kappa>0$,
$\tau_q\ge0$; no universal ordering of thermal and elastic speeds is assumed.

\textbf{Required parameters.} $C$, $e$, $\epsilon$, $\beta$, $p$, $\rho$, $c_E$,
$\kappa$, $\tau_q$, $T_0$ for both media, with source, temperature, units and a
consistent poling axis. If a reliable complete set is not available for a
material, choose a simpler documented pair. Do not synthesize missing constants.

\textbf{Characteristic problem and mode count (derive, do not assume).}
Substituting $\bm{y}=\bm{a}\exp\{i(k_x x + \xi z - \omega t)\}$ gives
$M(\xi,k_x,\omega)\bm{a}=0$ and $\det M=0$. The number of roots, their parity in
$\xi$, zero/gauge roots (for example a state with uniform $D_z$ or a constant
potential), and coalescences at $k_x=0$ are outputs of this derivation for the
chosen symmetry and electrical treatment. The count of modes admissible in each
region is the number of roots, computed from the derived $M$ for the actual
$(k_x,\omega)$, that satisfy the radiation or decay selection; it is not fixed in
advance. A preliminary exploratory computation with non-physical toy coefficients
(\textsc{executed}, outside this blueprint's evidence) suggested an even
polynomial and a $\xi=0$ root at $k_x=0$; this is a prompt to look for such
structure and is \emph{not} evidence for the physical system. Selection of
outgoing and decaying roots must use the flux and decay criteria of Section 5,
with continuation in $\omega$ and $k_x$.

\section{Contact conditions and configurations}

Take $z=0$ as a face. Define $\bm{t}=(\sigma_{xz},\sigma_{zz})^{\mathsf T}$,
velocity $\bm{v}=\dot{\bm{u}}$, electric normal displacement $D_z$ and heat flux
$q_z$.

\textbf{Primary mechanical contact (passive, diagonal).}
\begin{equation}
  \bm{t}^{(1)} = \bm{t}^{(2)} = \bm{t}_c, \qquad
  \bm{t}_c = K\jj\bm{u}\kk + C\jj\dot{\bm{u}}\kk, \qquad
  K\succeq0,\ C\succeq0, \tag{6}
\end{equation}
with $\jj\bm{u}\kk = \bm{u}^{(2)}-\bm{u}^{(1)}$ and the normal sign fixed so that
the interface-averaged dissipation equals the nonnegative form
$\tfrac12\jj\bm{v}\kk^{\mathsf H}C\jj\bm{v}\kk$ (\textsc{open} until verified in
gate G5b).

\textbf{Electrical contact.} Choose one: continuous potential benchmark, or
charge balance $D_z^{(2)}-D_z^{(1)}=\sigma_{\text{free}}$ with a passive
capacitive gap $D_z = C_s\jj\varphi\kk$ ($C_s\ge0$ in $\mathrm{F\,m^{-2}}$).
State whether free charge, storage or interfacial current is present.

\textbf{Thermal contact.} Mechanical coherence does not imply temperature
continuity:
\begin{equation}
  q_z^{(1)} = q_z^{(2)} = q_c, \qquad
  q_c = H_T\big(\vartheta^{(1)}-\vartheta^{(2)}\big), \qquad H_T\ge0, \tag{7}
\end{equation}
with the sign checked against the flux direction and entropy production. A
candidate time-averaged dissipation is $\tfrac12 H_T\lvert\jj\vartheta\kk\rvert^2/T_0$
(units $\mathrm{W\,m^{-2}}$); it must be derived, not assumed (G5b). No
off-diagonal $4\times4$ contact coefficients without an interface potential and
dissipation law.

\textbf{Equation-count requirement (derive per configuration).} For each
configuration the number of independent interface relations must equal the number
of unknown amplitudes after gauge removal; the rank is computed, not postulated.
Balance relations and jump relations together define the system. For the
interphase, counting is done at \emph{both} faces and includes internal
amplitudes.

\section{Scattering, energy partition and the scalar
reactive-versus-dissipative proposition}

\subsection{Scattering system and energy definitions}

For one incident mode $j$, assemble $B(\omega,k_x,\bm{p})\bm{a}_{\text{unk}}=\bm{b}_j$,
with $\bm{p}$ the sourced contact and layer parameters; solve by rank-aware linear
algebra, not raw determinants. Cramer ratios $r_m=N_m/\det B$ are analytical
descriptors only; writing $N_m=0$ is not a new theorem.

Define $R$, $T$, $A$ from a total-field flux derived for the adopted theory (gate
G4), evaluated on planes in the homogeneous half-spaces; evanescent and thermal
modes can produce cross terms, so channel-wise sums are used only where their
validity has been shown. The balance to be \emph{derived} and then verified is
\begin{equation}
  P_{\text{inc}} = P_{\text{refl}} + P_{\text{transm}} + P_{\text{bulk}}
  + P_{\text{contact}}, \tag{8}
\end{equation}
where $P_{\text{bulk}}$ is dissipation in the interphase (and in the half-spaces
between the evaluation planes and the contact), and $P_{\text{contact}}$ the
dissipation on the contact faces. The identity $R+T+A=1$ is \emph{not} a premise.

\subsection{Scalar model: assumptions}

These apply to Proposition 1 only.

\noindent (S1) One propagating channel on each side, normal incidence,
time-harmonic $e^{-i\omega t}$, lossless homogeneous media with real positive
impedances $Z_1, Z_2$.

\noindent (S2) Use the force-like variable $f=-\sigma_{zz}$ and velocity $v$. A
right-going wave obeys $f=Zv$, a left-going wave $f=-Zv$, and the time-averaged
power flux in $+z$ is $\tfrac12\operatorname{Re}(f\bar v)$. Write $Y_j=1/Z_j$,
$\Delta Y=Y_1-Y_2$, $\Sigma Y=Y_1+Y_2$, $\Sigma Z=Z_1+Z_2$.

\noindent (S3) The contact is zero-thickness, linear, local, of one of two
topologies, and passive. \emph{Topology V} (velocity jump): $f$ continuous,
$\jj v\kk = v_2-v_1 = -Y_d f$. \emph{Topology S} (force jump): $v$ continuous,
$\jj f\kk = f_2-f_1 = -Z_d v$. Absorbed power is
$-\tfrac12\operatorname{Re}(f\jj v\kk) = \tfrac12\operatorname{Re}(Y_d)\lvert f\rvert^2$
(V) and $\tfrac12\operatorname{Re}(Z_d)\lvert v\rvert^2$ (S); passive means
$\operatorname{Re}Y_d\ge0$, respectively $\operatorname{Re}Z_d\ge0$.

\noindent (S4) Combined two-port contacts, multimode problems and coupled
thermo-electric channels are outside the scope of the proposition.

The welded baseline is $Y_d=0$ (V) or $Z_d=0$ (S); then
$R_w = R_0 = \big((Z_1-Z_2)/(Z_1+Z_2)\big)^2$ and
$T_w = 1-R_0 = 4Z_1Z_2/(\Sigma Z)^2$.

\subsection{Derivation of the coefficients}

Medium 1: $f=f_i(1+r)$, $v=(f_i/Z_1)(1-r)$. Medium 2: $f=f_i t$, $v=f_i t/Z_2$.

\textbf{Topology V.} Force continuity gives $t=1+r$. The jump law
$t/Z_2 - (1-r)/Z_1 = -Y_d t$ becomes $(1+r)(Y_2+Y_d)=(1-r)Y_1$, hence
\begin{equation}
  r = \frac{\Delta Y - Y_d}{\Sigma Y + Y_d}, \qquad
  t = \frac{2Y_1}{\Sigma Y + Y_d}, \qquad
  R = \lvert r\rvert^2, \qquad
  T = \lvert t\rvert^2\frac{Z_1}{Z_2} = \frac{4Y_1Y_2}{\lvert\Sigma Y+Y_d\rvert^2},
  \qquad
  A = \frac{4Y_1\operatorname{Re}Y_d}{\lvert\Sigma Y+Y_d\rvert^2}. \tag{9}
\end{equation}

\textbf{Topology S.} Velocity continuity gives $t=Z_2(1-r)/Z_1$. The jump law
$\jj f\kk=-Z_d v$ gives $Z_2(1-r)-Z_1(1+r)=-Z_d(1-r)$, hence
\begin{equation}
  r = \frac{Z_2-Z_1+Z_d}{\Sigma Z + Z_d}, \qquad
  t = \frac{2Z_2}{\Sigma Z+Z_d}, \qquad
  R = \lvert r\rvert^2, \qquad
  T = \frac{4Z_1Z_2}{\lvert\Sigma Z+Z_d\rvert^2}, \qquad
  A = \frac{4Z_1\operatorname{Re}Z_d}{\lvert\Sigma Z+Z_d\rvert^2}. \tag{10}
\end{equation}

\textbf{Energy bookkeeping (\textsc{derived}).} For V, with $a=\operatorname{Re}Y_d$,
$\lvert\Sigma Y+Y_d\rvert^2-\lvert\Delta Y-Y_d\rvert^2 = 4Y_1Y_2+4Y_1a$, so
$R+T+A=1$ exactly; S is analogous.

\subsection*{Proposition 1}
Under (S1)--(S4):

\noindent\textbf{(a) Reactive contacts cannot suppress reflection.} If $Y_d=iB$
(V) or $Z_d=iX$ (S) with $B,X$ real, then
$R = \dfrac{\Delta Y^2+B^2}{\Sigma Y^2+B^2}$
(respectively $\dfrac{\Delta Z^2+X^2}{\Sigma Z^2+X^2}$) $\ge R_0$, with equality
only at $B=0$ (resp.\ $X=0$).

\noindent\textbf{(b) Exact suppression requires dissipation of the right
topology, and costs transmission.} $R=0$ with a passive contact is possible iff:
topology V with $Y_d=Y_1-Y_2$, which is passive iff $Z_1<Z_2$; or topology S with
$Z_d=Z_1-Z_2$, which is passive iff $Z_1>Z_2$. In both cases
$T=Z_{\min}/Z_{\max}$ and $A=1-Z_{\min}/Z_{\max}$. For a Voigt contact
$\bm{t}_c=K\jj\bm{u}\kk+C\jj\dot{\bm{u}}\kk$ in topology V one has
$Y_d=1/(C+iK/\omega)$, so exact $R=0$ requires $K=0$ and
$C=Z_1Z_2/(Z_2-Z_1)$ with $Z_1<Z_2$.

\noindent\textbf{(c) No passive zero-thickness contact increases transmission over
the welded baseline.} For every passive $Y_d$ (V) or $Z_d$ (S), $T\le T_w$, with
equality iff the contact vanishes. Consequently, in this model, no passive
contact achieves $R<R_w$ together with $T\ge T_w$.

\subsection*{Proof}

\noindent\textbf{(a)} For V, $r=(\Delta Y-iB)/(\Sigma Y+iB)$, so
$\lvert r\rvert^2 = (\Delta Y^2+B^2)/(\Sigma Y^2+B^2)$. Since $\Delta Y^2<\Sigma Y^2$
when $Y_1,Y_2>0$, adding the same $B^2$ to numerator and denominator of a fraction
below one increases it; equality requires $B=0$. S is identical with $X$ and $Z$.

\noindent\textbf{(b)} $R=0$ forces the numerator in (9) or (10) to vanish:
$Y_d=\Delta Y$ or $Z_d=Z_1-Z_2$. Passivity requires these to be real and
nonnegative, giving the stated orderings. Substituting, V gives
$T=4Y_1Y_2/(2Y_1)^2=Y_2/Y_1=Z_1/Z_2$ and S gives
$4Z_1Z_2/(2Z_1)^2=Z_2/Z_1$; $A=1-T$ by energy conservation. For the Voigt law with
$e^{-i\omega t}$ and $f=-\sigma_{zz}$, $-\bm{f}=K\jj\bm{u}\kk+C\jj\dot{\bm{u}}\kk$
and $\jj\bm{u}\kk=\jj\bm{v}\kk/(-i\omega)$ give
$\jj\bm{v}\kk=-\bm{f}/(C+iK/\omega)$, i.e.\ $Y_d=1/(C+iK/\omega)$, whose imaginary
part is zero iff $K=0$.

\noindent\textbf{(c)} For V, $\lvert\Sigma Y+Y_d\rvert^2=(\Sigma Y+a)^2+b^2$ with
$a=\operatorname{Re}Y_d\ge0$, $b=\operatorname{Im}Y_d$, which is $\ge\Sigma Y^2$
with equality iff $a=b=0$. Hence
$T=4Y_1Y_2/\lvert\Sigma Y+Y_d\rvert^2\le 4Y_1Y_2/\Sigma Y^2=T_w$. The argument for
S is the same with $\Sigma Z$. If $T\ge T_w$ then $Y_d=0$ (resp.\ $Z_d=0$), so
$R=R_w$.

\medskip
\noindent\textbf{Status.} The proposition and derivation are \textsc{derived}.
They are transmission-line algebra and are not presented as a new theorem.
Executed support (\textsc{executed}, Python): with $Z_2/Z_1=2$, minimization over
reactive elements returned $R_{\min}=0.1111=R_0$ in both orderings; resistive
matching returned $R=0$, $T=0.5000$, $A=0.5000$ with $R+T+A-1=0$ to machine
precision; random sampling of 2{,}000{,}000 passive contacts over five impedance
pairs (including $Z_1=Z_2$) and both topologies gave $\max(T-T_w)=0$, consistent
with (c); the closed-form $T$ matched the coefficient-based value to $10^{-16}$.

\textbf{Interpretation for the research.} Suppression by a zero-thickness passive
contact in this model is either \emph{impossible} (reactive, part a) or
\emph{dissipative and transmission-reducing} (parts b, c). When $Z_1\ne Z_2$,
write $\rho=Z_{\min}/Z_{\max}<1$. The transmitted fraction at exact matching is
$T=\rho$, while $T_w=4\rho/(1+\rho)^2$, so $T/T_w=(1+\rho)^2/4<1$
(\textsc{derived}). The proposition does \emph{not} decide the coupled,
multimode, oblique problem (S4).

\subsection{Scalar finite layer: benchmark and motivation for the primary
configuration}

For a lossless homogeneous layer of impedance $Z_\ell$ and phase $\phi=kd$
between the same media (S1), the input impedance is
$Z_{\text{in}} = Z_\ell\dfrac{Z_2+iZ_\ell\tan\phi}{Z_\ell+iZ_2\tan\phi}$ (the sign
of $i$ depends on the time convention; $R$ is unaffected). At $\phi=\pi/2$,
$Z_{\text{in}}=Z_\ell^2/Z_2$, so $R=0$ iff $Z_\ell^2=Z_1Z_2$, with $T=1>T_w$.
\textsc{derived}.

\textsc{executed}: for $Z_1=1$, $Z_2=2$, $Z_\ell=\sqrt2$, $\phi=\pi/2-10^{-9}$
gave $R=1.3\times10^{-19}$, $T=1$; at $\phi=0.3\pi/2$ the same layer gave
$R=0.0903<R_0=0.1111$ and $T=0.9097>T_w=0.8889$. This is the standard quarter-wave
result; it shows only that a finite lossless layer escapes the bound (c), which is
the motivation for H2, not a coupled-problem result.

\subsection{Coupled problem: criteria (hypotheses)}

For an incident wave with propagating reflected set $\mathcal{P}_{\text{refl}}$,
$R=0$ requires the amplitude of every channel in $\mathcal{P}_{\text{refl}}$ to
vanish. Each complex amplitude is two real conditions, so the number of free real
parameters needed generically depends on $\lvert\mathcal{P}_{\text{refl}}\rvert$
(\textsc{hypothesis}, a dimension count, not a transversality proof). Bulk LS
dissipation may make the minimum of $R$ nonzero; that outcome is reportable. Seek
an explicit relationship among contact compliance/damping, modal impedances,
incidence angle and a selected reflected amplitude; validate it first in reduced
limits (normal incidence, uncoupled elastic, weak coupling). If no closed form
survives, report the implicit condition, its sensitivity derivatives and the
mechanism. A determinant zero alone is not a new theorem.

\section{Finite interphase: primary configuration}

The interphase is a homogeneous piezothermoelastic layer $0<z<h$ obeying the same
field equations, with the passive contact laws of Section 4 on $z=0$ and $z=h$.
Its propagation phase is the additional parameter motivating H2. Both comparison
references are retained: the welded two-half-space case and the
imperfect-contact two-half-space case. All use the same sourced outer media,
incident mode, frequency-angle convention, and flux normalization. A finite
interphase must be compared with both, and its added free parameters disclosed.

\textbf{Unknown-amplitude and equation bookkeeping (derive for the actual
system).} Count the admissible reflected, interphase (both directions) and
transmitted amplitudes from the derived $M$ for each region and the incident mode;
count independent relations at both faces; verify rank after removing gauge or
null modes. Neither the unknown total nor the number of relations is fixed in this
blueprint. Use a scaled scattering or impedance-type formulation rather than a
naive transfer matrix when attenuating modes grow exponentially.

\textbf{The $h\to0$ reduction (to be derived, not assumed).} Write the face laws
and the layer as maps on the state $\bm{s}=(\bm{q},\bm{F})$, with
$\bm{q}=(u_x,u_z,\varphi,\vartheta)$ and generalized fluxes
$\bm{F}=(\sigma_{xz},\sigma_{zz},D_z,q_z)$. If the layer admits
$d\bm{s}/dz=\bm{A}\bm{s}$, then the composite map is
$\bm{s}(h^{+}) = \bm{P}_2 e^{h\bm{A}}\bm{P}_1\bm{s}(0^{-})$ with $\bm{P}_i$ the
face maps. The $h\to0$ limit of this composition must be \emph{computed},
including whether it yields a jump law of the same form (flux continuity with a
$\bm{q}$ jump), and what the first-order dependence on $h$ is. The expected form is
an effective compliance consisting of the two face compliances plus an $O(h)$
contribution from the layer itself, with sign and frame conventions fixed; this is
a \textsc{hypothesis} to be confirmed or corrected by the derivation. Two
compliant faces do not in general give a welded contact, and no equivalence with
the baseline is assumed. Independent limits to test: ideal faces ($K\to\infty$)
with identical media (no reflection), and the scalar matched layer of
Section 5.4.

\textbf{Switch rule.} If, with sourced parameters, the interphase offers neither a
transmission-preserving reflection reduction nor a defensible analytic advantage
over the baseline (H2 rejected), the paper is reported with the baseline as the
main configuration and interphase-focused claims and title removed.

\section{Mandatory validation protocol}

Validation is a compulsory scientific result. No statement of agreement is
permitted without an aligned quantity, a stated source and a difference. Nothing
below is marked completed by this blueprint. An energy residual is an internal
consistency check and is not independent validation.

\subsection{Validation route V1: published compatible reflection/transmission
results (conditional)}

No single published paper is mandated. A compatible independent benchmark is
preferred over an author-group self-benchmark; when inaccessible, independent
analytical and code-based cross-checks are required, and the published-comparison
gate is reported as unavailable rather than falsely passed. A published result may
be used only after a documented \emph{compatibility check} (a file
\texttt{compat\_check.md} per candidate) establishing, from the full accessible
text: bulk constitutive model and thermal theory; the set of modes retained on
each side; boundary-condition type and sign conventions; material symmetry and
poling direction; incident wave, frequency or angle definition and normalization;
and which of the paper's parameters can be switched off to reach the present
model. Candidates:

\begin{itemize}
  \item Barak et al.\ \texttt{10.1007/s12648-022-02568-w}: its model includes
    nonlocality, memory-dependent derivatives and a hyperbolic two-temperature
    option with dual-phase lag, and its first medium is elastic with two reflected
    and four transmitted waves (abstract). A usable limit requires nonlocality,
    memory and two-temperature effects to be switched off and the thermal law to
    reduce to LS as shown in the paper, and the heterogeneous mode count of the
    elastic/piezothermoelastic problem must be derived separately from the
    symmetric one. If the full text, constants or normalization are unavailable,
    or the limit cannot be demonstrated, mark the candidate \textbf{NOT USABLE AS
    BENCHMARK} and use another route.
  \item Guha, Singh and Das \texttt{10.1080/17455030.2019.1610198} and Karmakar,
    Sahu and Goyal \texttt{10.1177/10812865221099502} for imperfect contact
    between piezothermoelastic half-spaces, under the same compatibility test
    (including their contact-law definitions).
\end{itemize}

When a source is only partly compatible, report a \emph{related-physics
comparison}, not a validation. Numbers are taken from tables or data when
available; digitized values carry a stated extraction uncertainty. Discrepancy is
$\delta_i = \lvert Q_{\text{calc}_i} - Q_{\text{ref}_i}\rvert$, with a relative
error only where $\lvert Q_{\text{ref}_i}\rvert$ is appreciably nonzero.

\subsection{Validation route V2: exact analytical limiting cases (independent of
V1)}

Run each as a separate test with an explicit parameter path:
\begin{enumerate}
  \item Identical welded half-spaces: zero reflected field for an incident bulk
    mode (subject to the declared gauge and mode basis).
  \item Uncoupled elastic normal incidence: $r=(Z_2-Z_1)/(Z_2+Z_1)$ up to the
    amplitude convention, $R=\lvert r\rvert^2$, $T=1-R$ for lossless media with
    $Z=\rho c_L$.
  \item Classical elastic P--SV oblique incidence against an independently
    formulated welded and spring--dashpot solution, with critical-angle and
    evanescent treatment.
  \item Scalar proposition limits: with couplings off, the contact reproduces
    Eqs.\ (9)--(10), the bound $T\le T_w$ and the exact-suppression values of
    Proposition 1; the interphase reproduces the quarter-wave result of
    Section 5.4.
  \item Perfect-versus-imperfect contact along stated paths (high stiffness,
    vanishing compliance and vanishing damping are distinguished).
  \item Thermal and coupling reductions: $\tau_q\to0$ to the Fourier limit;
    $\beta,p,e\to0$ to uncoupled fields; these are not the same as deleting a
    field.
  \item Interphase: the independently derived $h\to0$ reduction of Section 6.
\end{enumerate}

\textbf{Preliminary exploratory check (outside this blueprint's evidence).} In the
audit session a classical elastic P--SV spring--dashpot solver on dimensionless toy
media was run (\textsc{executed}): energy residual $\le6\times10^{-16}$ in 12 cases
and the welded normal-incidence $R=0.033630$ matched the Fresnel value. This does
not validate any piezothermoelastic solver and is not cited as validation.

\subsection{Validation route V3: mathematical and numerical correctness}

\begin{enumerate}
  \item Eigenmode residual $\lVert M\bm{a}\rVert/(\lVert M\rVert\lVert\bm{a}\rVert)$
    with an independent root formulation (polynomial versus state-space/boundary
    matrix); outgoing selection by flux and decay, not by the sign of
    $\operatorname{Re}\xi$ alone.
  \item Boundary residual and rank: normalized $\lVert B\bm{a}-\bm{b}\rVert$,
    singular values, condition number, independent-equation count including gauge
    constraints, near-degenerate roots and critical angles.
  \item Convergence: at least three tightened root/eigen tolerances and two
    numerical resolutions where relevant; scattering-matrix versus direct
    linear-system agreement; record changes in $R$, $T$, roots and the location
    of any minimum.
  \item Robustness against artificial minima: coarse independent grid with local
    refinement; repeat near candidate minima at higher precision; exclude branch
    swaps, nonphysical parameters, wrong attenuation signs and absorption
    mislabeled as transmission.
\end{enumerate}

\subsection{Compulsory energy gates}

\textbf{E1 LS energy-flux derivation.} From the adopted constitutive relations,
LS conduction and Maxwell (quasi-static) equations, derive the coupled
electromechanical--thermal energy identity, the normal total-field flux and its
thermal term, and the stored-energy/entropy expressions including the
$\tau_q$-dependent terms. Verify positivity of bulk dissipation under the
stability conditions. No flux expression is adopted from another thermal theory.

\textbf{E2 Interface dissipation.} Derive the mechanical, thermal and electrical
(if any) contact dissipation as nonnegative forms from the contact laws, including
sign conventions, and cross-check them against the flux difference across the
contact.

\textbf{E3 Numerical energy balance.} Evaluate Eq.\ (8) with the derived flux and
dissipations on evaluation planes in the half-spaces; require convergent residuals
for lossless reductions, and a physically signed, convergent balance for
dissipative cases.

\textbf{Failure of any energy gate blocks the reporting of $R$, $T$ or $A$.}

\subsection{Evidence acceptance and presentation}

Produce one main-text validation figure (a compatible published overlay or an
analytical limit) and one validation table with reference or analytical value,
assumptions, present value, absolute and relative discrepancy, solver settings and
status. Numerical tolerances are fixed \emph{before} reporting; values such as
$10^{-6}$ for benchmark discrepancy or $10^{-8}$ for normalized residuals are
targets to be tested, not guarantees. Where the reference precision or conditioning
is weaker, report the limitation.

\begin{center}
\begin{tabular}{@{}p{0.20\textwidth}p{0.48\textwidth}p{0.16\textwidth}@{}}
\toprule
\textbf{Evidence test} & \textbf{Proposed acceptance rule (justify against
scales)} & \textbf{Status} \\
\midrule
V1 published comparison & Compatibility check passed; discrepancy within source
  precision and demonstrated convergence & Pending \\
V2 exact limits & Analytical relation reproduced to tight tolerance, including
  the scalar proposition & Pending \\
E1--E3 energy gates & Derived flux, signed dissipations, convergent balance &
  Pending \\
Interface rank and residual & Independent rank correct, small normalized
  residual, no spurious gauge amplitudes & Pending \\
Numerical minimum & Persists under refinement and perturbation; $\Delta T$ and
  $A$ reported & Pending \\
\bottomrule
\end{tabular}
\end{center}

\textbf{Stop rule.} If any mandatory validation fails, no submission-ready results
are generated; repair the model or narrow the hypothesis. A genuine null result
may be reported but does not replace validation.

\section{Staged gates and honest failure routes}

\begin{center}
\begin{tabular}{@{}p{0.07\textwidth}p{0.50\textwidth}p{0.30\textwidth}@{}}
\toprule
\textbf{Gate} & \textbf{Required evidence} & \textbf{Decision on failure} \\
\midrule
G0 & Full-text prior-art comparison: Guha; Karmakar; Gupta-group baselines;
  Wang; Ghazaryan; at least one recent layered comparator. & Refine one
  mechanism, not the field. \\
G1 & Sourced complete bulk constants, dimensional check, symmetry,
  thermodynamic stability, thermal parameters. & Replace pair or simplify
  theory. \\
G2 & Derived $M$, root count and parity for the chosen symmetry, gauge/zero
  roots, outward selection, convergence. & Correct model and boundary closure. \\
G3 & Validation routes V1 (conditional) and V2 (exact limits) completed,
  including the $h\to0$ derivation and the scalar proposition limits. & Repair
  algebra before parametric plots. \\
G4 & E1: derived LS energy identity and flux. & Do not quote any energy ratio. \\
G5 & E2 and E3: derived interface dissipation and convergent numerical energy
  balance. & Do not quote any energy ratio. \\
G6 & Parameter search within the pre-declared box and thresholds
  $\varepsilon_R$, $\delta_T$ reports H1 and H2 outcomes with $\Delta R$,
  $\Delta T$, $A$. & Report the negative result; adjust only through a logged
  decision. \\
G7 & Competitor comparison shows a specific mathematical or physical distinction
  with substantiated benefit. & Retitle as a narrower study or stop. \\
\bottomrule
\end{tabular}
\end{center}

No guessed curves, fabricated supplementary data, unsupported novelty claim, or
numerical result without code, constants, tolerances and raw outputs. A compiled
LaTeX file is not validated physics.

\section{Flexible workflow and permitted adjustments}

If an anti-reflection hypothesis fails, the permitted responses are: change
contact parameters within physically justified and sourced ranges; change the
interphase thickness or material; change the material pair; change the incident
mode or angle range; replace the contact law with another passive law. Each
response is logged. The following are \emph{not} permitted: fabricating or
adjusting results; tuning a parameter outside its justified range solely to
produce a minimum; changing $\varepsilon_R$, $\delta_T$, the parameter box or the
validation criteria after seeing results without a logged justification that is
reported in the paper; reporting dissipation-driven suppression as
useful-transmission anti-reflection. Parameter optimization, if used, is confined
to the declared box and its settings and outputs are retained.

\section{Proposed paper organization and evidence plan}

\begin{enumerate}
  \item Introduction: motivation, competitors, precisely limited question.
  \item Formulation: geometry, symmetry, three coupled categories, LS theory.
  \item Contact and interphase conditions: spring--dashpot, electrothermal
    closure, passivity, equation count.
  \item Wave solutions: characteristic roots, admissibility, scattering algebra,
    interphase reduction.
  \item Reflection suppression and power: scalar proposition, coupled criteria,
    flux balance, loss--transmission separation.
  \item Validation and results: validation figure and table first, convergence,
    baseline-versus-interphase comparison, sensitivity, limitations.
  \item Discussion and conclusions: what was shown, what was not, scope for
    future extensions.
\end{enumerate}
Planned figures (earned by findings): schematic; root structure; $R,T,A$ versus
angle or frequency; contact-parameter response; baseline versus interphase
comparison; energy residual and convergence. Target 5--7 figures and 2--4 tables,
not a fixed count.

\section{Decision record (mandatory)}

For each change to an assumption, mechanism, parameter box, threshold or model,
add a dated row: original assumption; reason; calculation or citation; revised
model; effect on boundary conditions and flux accounting; retests; effect on title
and novelty; status. Negative results are retained. Not fixed: material pair,
$\tau_q$, polarization, contact parameters, $\varepsilon_R$, $\delta_T$ (until
declared), final title. Fixed unless the whole first paper is reconsidered:
piezothermoelastic wave dynamics, scattering analysis, coupled-field consistency,
and an independently validated energy account.

\section{Referee checklist}

A reviewer should cite section, quote the claim, classify
(critical/major/minor/preference), give the calculation or reference, state
assumptions and the smallest correction. Specifically check: LS signs and
pyroelectric term; mode count and gauge roots from the derived $M$; contact
passivity and equation rank; the $h\to0$ derivation; the scalar proposition
(assumptions S1--S4) and whether H1 as stated is consistent with it; the energy
gates; benchmark compatibility; full-text overlap rather than abstract-level
originality.

\section{Final research management and agent work orders}

\textbf{Ownership.} The principal investigator/supervising research team retains
control of the final physical model, material pair, hypotheses, manuscript
structure, and approval. Neither Claude nor Arena may introduce a new physical
field, replace piezothermoelasticity, switch to device engineering, or promote a
speculative finding into a conclusion. Scientific corrections are welcome;
uncontrolled programme redesign is not.

\subsection*{Single-paper scope lock}
\begin{itemize}
  \item \textbf{Primary paper:} passive reflection suppression of an incident
    quasi-longitudinal coupled wave in two outer piezothermoelastic media with a
    finite piezothermoelastic interphase and generalized passive imperfect
    contact. Oblique incidence is investigated after the normal-incidence limit is
    stable.
  \item \textbf{Compulsory comparators:} welded two-half-space model and
    imperfect-contact two-half-space model, with like-for-like material parameters
    and clear parameter counts.
  \item \textbf{Optional within-paper decisions:} material selection, one
    physically supported contact law, and interphase thickness. Switch to
    two-half-space primary formulation only if the predeclared gate in Section 6
    fails. Such switching is disclosed, not hidden.
  \item \textbf{Deferred:} Gurtin--Murdoch surface elasticity,
    Klein--Gordon/NSGT nonlocality, additional coupled fields, experiments and
    device fabrication. These are not required to rescue a disappointing result.
\end{itemize}

\subsection*{Phase ownership and deliverables}
\begin{center}
\begin{tabular}{@{}p{0.20\textwidth}p{0.36\textwidth}p{0.38\textwidth}@{}}
\toprule
\textbf{Phase} & \textbf{Primary assignment} & \textbf{Required acceptance
evidence} \\
\midrule
P0 Claude: prior-art and constants & DOI/title/author cross-check; closest
  full-text comparison when accessible; complete sourced constitutive parameter
  table; explicit missing fields. & \\
P1 Claude: constitutive/boundary derivation & Thermodynamic potential, corrected
  LS coupling, full contact power balance, independent BC rank and dimension
  checks. & \\
P2 Arena: symbolic and numerical solver & Reproducible scripts, eigenmode
  residuals, root selection, direct boundary solve and stable scattering
  calculation, raw arrays and environment details. & Reproducible scripts,
  eigenmode residuals, root selection, direct boundary solve and stable
  scattering calculation, raw arrays and environment details. \\
P3 Arena: independent validation & Classical elastic limits, identical-media,
  zero-thickness composition, scalar proposition, published benchmark when
  compatible, solver cross-check, converged energy residuals. & \\
P4 Claude: critical independent audit & Examine Arena outputs and scripts,
  identify genuine errors with equations/line numbers, and distinguish verified
  defects from preferences. & \\
P5 Supervisor: final interpretation & Approve or narrow hypotheses, choose final
  title, ensure novelty evidence, approve manuscript only when all critical gates
  pass. & \\
\bottomrule
\end{tabular}
\end{center}

\subsection*{Mandatory evidence contract for every agent}
Each deliverable must state (i) which source files were actually opened,
(ii) commands/code run versus not run, (iii) exact input constants and provenance,
(iv) equations and units, (v) residuals/tolerances, (vi) deviations from this
working blueprint, and (vii) unresolved or failed tests. Labels \textsc{derived},
\textsc{executed}, \textsc{verified source}, and \textsc{unverified} must be used
faithfully. Do not write ``validated'', ``Q1 ready'', ``supplementary files
provided'', or ``simulation completed'' without matching evidence. Agent
disagreement is resolved by reproducible calculation, not by a vote or a numerical
score.

\subsection*{Stage-gate decisions}
Gate A: complete sourced constants and well-posed field/interface model. Gate B:
independent analytical limits and stable modal/scattering solver. Gate C: derived
coupled power accounting and converged passive energy balance. Gate D: a
meaningful, independently validated reflection--transmission finding relative to
both baselines. Gate E: source-grounded novelty and clear mathematical
contribution. A failed gate stops the associated claim, not necessarily the full
research programme. No manuscript result may be presented as experimental or
computed without passing its relevant gates.

\subsection*{Additional mathematical closure checks before numerical claims}
\begin{enumerate}
  \item The scalar Proposition 1 remains a \emph{benchmark theorem} for the stated
    scalar topologies, not a universal no-go theorem for coupled
    piezothermoelastic scattering. Its traction sign and Voigt mapping must be
    checked independently against the adopted global traction convention.
  \item The interface electrical charge-storage relation needs a complete
    charge-balance and time-harmonic current formulation. Never impose both
    arbitrary free surface charge and independent potential/$D$ conditions
    without counting constraints.
  \item The proposed thermal-contact dissipation must be derived with consistent
    absolute-temperature factors. Check both dimensions and nonnegative entropy
    production rather than adopting a quadratic expression by analogy.
  \item A dissipative half-space has complex modes and possible interference
    terms. The definition of reflected/transmitted power, evaluation planes, and
    separation from bulk loss must be fixed before quoting $R$, $T$, or $A$.
  \item Hypothesis H2 requires comparisons with both a welded two-half-space
    reference and the imperfect-contact reference; a superior result obtained
    using more tunable parameters is not by itself proof of a new physical
    mechanism.
\end{enumerate}

\section{Disposition}

\textbf{Working paper identity:} Reflection suppression of coupled waves in
piezothermoelastic media through generalized imperfect interface conditions.

\textbf{Primary candidate configuration:} finite piezothermoelastic matching
interphase with passive spring--dashpot faces. \textbf{Baseline:} two-half-space
generalized imperfect interface.

\textbf{Approval status:} final working blueprint authorized for scoped
feasibility and validation; not an approval of novelty claims or manuscript
submission. No coupled-field numerical result has been produced.

\appendix
\section{Historical change log: V4.1 $\to$ V4.2}

\begin{center}
\begin{tabular}{@{}p{0.04\textwidth}p{0.22\textwidth}p{0.64\textwidth}@{}}
\toprule
\# & Location & Change \\
\midrule
1 & Section 5 (new 5.2--5.4) & Added scalar assumptions (S1--S4), derivations of
  $r,t,R,T,A$ for both contact topologies, Proposition (a)--(c) with proof,
  executed checks, and the interpretation of the transmission cost; added the
  scalar lossless-layer benchmark. \\
2 & Section 2, H1 & H1 no longer presumes simultaneous reduction of $R$ with
  $T\ge T_w$. Now uses pre-declared $\varepsilon_R$, $\delta_T$, states the
  derived counterexample, conditions for possible coupled-case success, and
  falsification rules F1--F3. \\
3 & Sections 1, 2, 6 & Finite interphase is the primary candidate configuration
  (H2); two-half-space interface is the comparison baseline. V4.1 had the
  interphase as a conditional extension. \\
4 & Section 7, V1 & Exact reproduction of Barak et al.\ replaced by a conditional
  benchmark with a documented compatibility check (nonlocal, memory,
  two-temperature limits; heterogeneous mode count); fallback and ``not usable''
  status added. \\
5 & Section 7, V1 and V2 & Two independent validation routes stated: published
  compatible $R/T$ results \emph{and} exact analytical limits (including scalar
  proposition and quarter-wave limits). \\
6 & Section 7, E1--E3; Section 8 & LS energy-flux derivation, interface
  dissipation and numerical energy balance made compulsory gates blocking any
  energy-ratio claim; gates G4, G5 restructured. \\
7 & Sections 3, 6 & Mode count and $h\to0$ reduction are to be derived, not
  assumed: removed the pre-set counts; added the state-space composition for the
  reduction with the expected form marked as a hypothesis; toy-polynomial
  observation explicitly excluded as evidence. \\
8 & Section 9 & Added explicit permitted and prohibited adjustments consistent
  with the flexible workflow. \\
9 & Section 2 & Ghazaryan et al.\ DOI record checked (Crossref); competitors'
  read-status stated. \\
10 & Sections 1, 12 & Gurtin--Murdoch and Klein--Gordon excluded unless a
  demonstrated need is logged; disposition updated. \\
\bottomrule
\end{tabular}
\end{center}

\end{document}
